\chapter{Home}
\label{chap:home}

We have followed the framework from the timeless void (Chapter~\ref{chap:void})
through the first dawn (Chapter~\ref{chap:dawn}) down the recursion to its
children. This chapter arrives home---at our own universe---and finds that the
parts of the sky cosmology has had to \emph{add by hand}, the dark sector, are
here not additions at all but the plain look and feel of living inside a black
hole. Dark matter, dark energy, the microwave background, the great rings in the
galaxy distribution: read in this framework, each is a feature of our address, and
together they are a family portrait of the parent on the far side of our horizon.
And they carry a single deeper message, the one the book is named for: our
beginning has not ended. The textbook Big Bang is a finished event in the deep
past; ours is a surface still pressing on us from the edge of the sky, our parent
still feeding matter across it. The dark sector is exactly the part of that
\emph{ongoing} dawn we are able to measure---so far from being awkward add-ons,
dark matter and dark energy are the standing proof that we live not after our
dawn but inside it.

\section{We are a child}
\label{sec:we-are-child}

One observation settles our place in the family tree before any detailed
accounting. Our universe contains dark matter and dark energy, and in this
framework both \emph{require a parent}: dark matter is the parent's mass felt
through the persistent bounce membrane, and dark energy is the parent's ongoing
accretion (this chapter makes both precise). An original-generation universe,
nucleated straight from the void, has no parent and so would have neither. Their
presence in our sky is therefore a direct reading:
\[
  \text{(observed dark matter and dark energy)} \;\Longrightarrow\;
  \text{we are a child universe---a black-hole universe, generation } n\ge 1 .
\]
We are not a firstborn. The darkness of the sky is the evidence.

\section{The persistent channel}
\label{sec:persistent-channel}

The bounce membrane is not a one-time event but an ongoing interface
(Chapter~\ref{chap:bounce}). As long as the parent black hole keeps accreting,
matter keeps falling in, reaching the bounce density, and crossing into our
interior---throughout the parent's entire life. But the injection geometry is
peculiar, and everything in this chapter follows from it: new matter does not
appear at some location in our three-dimensional space. It appears \emph{at the
membrane}, which from inside is our cosmic horizon at maximum lookback---our
Big-Bang surface. Injected matter therefore manifests as gravitational influence
arriving from the cosmic boundary, not as particles materialising among the
galaxies.

\begin{quote}\small
\textbf{Why you can seed a child but never visit one.} It is worth seeing what
kind of surface this is, because it answers a natural question: could one, granted
every speculative tool relativity allows---warp drives, traversable wormholes,
exotic matter---\emph{travel} to the parent, or to a baby universe forming in our
sky? No, in both directions, for the same reason. The membrane is not a wall in
space, which a warp bubble could defeat, but a surface in \emph{time}: a bounce, a
low-entropy boundary. The parent lies behind our Big-Bang surface, in our
\emph{past}; you cannot drive to your own Big Bang. And a baby universe is a real
black hole you \emph{can} fall toward---but its time runs forward from its bounce,
so the horizon delivers you to its low-entropy first instant, not to any interior
``now.'' You could become a speck in its initial conditions; you can be an
ancestor, never a tourist. The membrane-as-Big-Bang structure is the framework's
own chronology protection: you may fall down the recursion, never climb back up.
\end{quote}

\section{Dark matter: the parent's gravity}
\label{sec:dark-matter}

The matter on the parent's side of the membrane has mass--energy, so it curves our
space; but it does not enter our interior, exchange photons with our matter, or
annihilate against it. That is precisely the phenomenology of dark matter:
halo-like mass around galaxies, extra rotational binding, lensing without visible
source, large-scale correlations---all gravitational, none luminous. In this
framework it is not a new particle but the live gravitational signature of the
parent's ongoing, biased accretion.

\paragraph{Gravity felt before light seen.} Because the membrane is our
\emph{past} horizon, the parent's mass is already in our metric---it curves our
space \emph{now}---while its light is still climbing in toward us from that
receding past surface, and will for cosmic ages. So the dark sector's
gravitational pull \emph{leads} its electromagnetic visibility: we feel matter we
cannot yet see (\Cref{fig:lookback}). ``Dark'' matter is, on this reading, exactly
gravity arriving ahead of light from the membrane.

\paragraph{Perhaps the missing antimatter.} If the bounce has any chirality-selective
character, the rejected handedness remains at the membrane---gravitationally
present, geometrically separated from our interior, unable to annihilate against
our matter. Cosmology's ``missing antimatter'' and its ``dark matter'' may then be
one phenomenon. The reading even carries a number: if a fraction $f$ of the
infalling matter passes the extrusion, the dark-to-baryonic ratio is about
$(1-f)/f$, so the observed $\sim 5\!:\!1$ implies $f\sim1/6$---a target a
first-principles Einstein--Cartan calculation should hit, and a sharp way to be
wrong (Chapter~\ref{chap:sims}).

\paragraph{What distinguishes it from a particle.} Geometric dark matter need not
clump exactly as cold particles do, need not appear in any direct-detection
experiment, and should trace the parent's projected mass rather than our own
luminous structure---tests taken up in Part~II (Chapter~\ref{chap:nowobs}).

\begin{figure}[t]
  \centering
  \includegraphics[width=\textwidth]{lookback}
  \caption{Gravity felt before light seen. The parent's matter sits at (and beyond)
    our past cosmological horizon---the membrane---so its mass is already in our
    metric and pulls on us now, while its light is still climbing in from that
    receding surface. The dark-matter gravitational signature therefore leads its
    electromagnetic visibility. Reproducible as
    \texttt{figures/scripts/ch06\_lookback.py}.}
  \label{fig:lookback}
\end{figure}

\section{Dark energy: the parent's growth}
\label{sec:dark-energy}

The parent is not a static reservoir; it is a growing black hole, feeding from its
own universe. As its mass climbs, the membrane's geometry relative to our interior
changes, and from inside that shows up as an extra outward push on our
expansion---the acceleration that standard cosmology attributes to dark energy.
This makes dark energy a \emph{rate}: to leading order its strength tracks the
\emph{time-derivative} of the parent's accretion. A parent still gorging sustains
or strengthens it; a parent whose surroundings are thinning, or which is turning
toward Hawking-dominated mass loss, yields weakening dark energy---and, in the far
future, a reversal, as the membrane shrinks and our expansion decelerates. So
dark energy here is neither constant nor eternal: it is a readout of our parent's
life story, and the recent hints from DESI that its equation of state may be
evolving are exactly where to confront it (Chapter~\ref{chap:nowobs}).

\section{The microwave background: the parent's glow}
\label{sec:cmb}

The membrane wears, we have seen, three faces: the parent's event horizon from
outside, our cosmic horizon at maximum lookback from inside, the Big-Bang surface
in the standard telling. A black hole's horizon radiates---Hawking's
result---and from our interior that radiation arrives from every direction, at
maximum cosmic lookback, with a thermal spectrum. That is the cosmic microwave
background. The parent's Hawking temperature at emission is essentially zero for so
large a hole, but the radiation has been redshifting through our entire expansion;
the $2.7\,\mathrm{K}$ we measure is what survives. The identification sits well
with what the background \emph{is}: a near-perfect Planck spectrum (Hawking
radiation is thermal), near-perfect isotropy (the membrane is smooth across the
parent's horizon), and anisotropies at the $10^{-5}$ level (membrane inhomogeneity
and a little parent rotation). The one genuine debt is the acoustic peak
structure, which standard recombination fits beautifully and which a bounce
cosmology must re-derive---the central calculational challenge, taken up in
Chapter~\ref{chap:sims}.

If the parent rotates, its axis imprints a preferred direction on the membrane,
and so on our earliest sky and the structure that grows from it. The reported
large-angle CMB alignments---the ``axis of evil''---and the JWST galaxy-spin
handedness asymmetry are candidate signatures of that single inherited axis; the
framework's sharp prediction is that they should \emph{share} it
(Chapter~\ref{chap:nowobs}).

\section{Lumpiness: the texture of the meal}
\label{sec:lumpiness}

What the parent feeds in is not smooth. A growing hole eats streams, clumps, the
cosmic web of its own universe, so the matter projected through our membrane
carries the parent's inhomogeneity---structure on scales unrelated to our own
internal history. This sharpens into a structural signature of the recursive
supraverse. Trace lumpiness down a lineage: the void is homogeneous; an OGU, with
no parent, inherits none and (lacking dark matter) builds only weak internal
structure---it is the minimal-lumpiness rung, and \emph{this is exactly what
$\Lambda$CDM assumes}, a smooth near-singular start with structure built only from
within. In the language of this framework, $\Lambda$CDM describes an OGU. But a
child inherits its parent's lumpiness \emph{and} adds its own, so large-scale power
\emph{accumulates} generation by generation toward a fixed point, with the biggest
jump at the first child, where a parent's lumpiness appears for the first time
(\Cref{fig:lumpiness}).

The observable falls straight out: our universe should show \emph{more}
ultra-large structure than any smooth-start model predicts. And it does---the
Giant Arc, the $\sim\!400\,\mathrm{Mpc}$ Big Ring~\citep{lopez2024bigring}, the
Sloan Great Wall, the KBC void, all exceeding the $\sim\!260\,\mathrm{Mpc}$
homogeneity scale and resisting assembly in the time available. In $\Lambda$CDM
these are an embarrassment; here they are simply the inherited rung. Lumpiness thus
joins dark matter and dark energy as a third, independent sign that we are a child
and not a firstborn. The model is illustrative---the intrinsic and transferred
contributions are not yet derived from the membrane projection---but its robust
content is a direction a singular start cannot produce: monotonic excess
large-scale power for every generation past the first.

\begin{figure}[t]
  \centering
  \includegraphics[width=\textwidth]{lumpiness}
  \caption{The lumpiness ladder. \emph{Left:} large-scale lumpiness climbs from the
    homogeneous void through the OGU (minimal---the $\Lambda$CDM smooth-start case)
    to the first child (us; the biggest jump, where the parent's lumpiness first
    appears) and on toward a fixed point. \emph{Right:} the smooth-start floor
    $\sigma(R)$, computed with CAMB, falls to $\sim0.01$ by the $260\,\mathrm{Mpc}$
    homogeneity scale, so the observed ultra-large structures (markers) sit far
    above what a singular start builds---the inherited rung that marks us a child.
    Reproducible as \texttt{figures/scripts/ch06\_lumpiness.py}.}
  \label{fig:lumpiness}
\end{figure}

\section{Where we sit}
\label{sec:where-we-sit}

It is worth saying our address in one breath, because it has a quiet joke in it.
The black holes we watch form in the sky---in galactic centres, in collapsing
stars---are, in this framework, the \emph{next} generation: each is the seed of a
child universe, a $\mathrm{BHU}_2$ to our $\mathrm{BHU}_1$. And the picture
cosmology drew for our \emph{own} beginning, a Big Bang from a near-singular,
smooth, parentless start, is the description of an \emph{OGU}---a first universe,
the generation \emph{above} us (Chapter~\ref{chap:dawn}). So we have spent a
century writing, with great care and success, the biography of our
grandparent---the parentless OGU at the root of our lineage---and taking it for
our own. We are not at the root, and not at the leaves. We sit in between: a child
watching its children form, having mistaken its grandparent's story for
autobiography. The dark sector is exactly the part of the story that gives the
mistake away---the parent that the OGU never had, pressing on us through the
membrane.

\medskip
\noindent None of this is bolted on. Dark matter, dark energy, the background
glow, the great rings: each is a different glimpse of the same parent pressing
through the same membrane, and together they place us---unspecially---inside one
black hole among uncountably many. That is the framework's picture of home.
Whether the picture is \emph{true} is the question Part~II takes up: what we can
already see, what the next decade will decide, and what only computation can yet
settle.
